\chapter{Data}\label{ch:data}
\begin{goals}
\textbf{In this chapter} you open measurements from any delimited text file, learn exactly how wuffle
decides what each column means, and work with a dataset's rows in the table.
\end{goals}

A dataset is where everything starts. wuffle reads the files geologists have: a spreadsheet
exported as CSV, a column of attitudes typed into a notebook, a log from another program with its own
column names. It never refuses one. What it can't read it counts, and what it had to guess it tells you.

\section{Opening a file}\label{sec:data-open}

Choose \menu{File>Open data or a project…} (\key{Ctrl}+\key{O}) or drop files on the window. Each
text file becomes a dataset named after the file, without its extension, and is drawn as a layer on the
active diagram. If the project has no diagram yet, wuffle makes \ui{Net 1} first. Opening a file is one
undo step, however many things it adds (\cref{sec:undo}).

Several files can be opened at once. A name already in the project gets a number: a second
\file{joints.csv} becomes \code{joints (2)}.

\section{How wuffle reads a file}\label{sec:data-read}

\subsection{Lines, delimiters and numbers}

\begin{itemize}
\item Blank lines are skipped, and so are comment lines starting with \code{\#} or \code{//}.
\item The \term{delimiter} is the one most lines agree on, out of comma, semicolon, tab and
\code{|}: the one that splits at least three lines in four into the same number of fields. One odd line
doesn't make the file unreadable, even a short one: in \code{dd,dip} / \code{10,20} / \code{hello} /
\code{30,40} the comma still wins, and \code{hello} is counted as skipped. When two delimiters split the
lines equally well, the one whose fields read as numbers wins. With none of them, wuffle splits on
spaces.
\item Fields in double quotes may hold the delimiter: \code{"open, iron-stained"} is one field.
\item A \term{decimal comma} is recognized when the comma isn't the delimiter: \code{120,5;40,2}
reads as 120.5 and 40.2, with or without a header line. (It splits evenly on the comma too, but there
\code{5;40} is no number, so the semicolon wins.)
\item The first line is a \term{header} when it holds a word, something that isn't a number or an
attitude. Without one, the columns are called \col{c1}, \col{c2}, and so on.
\item A line that can't be split into at least two fields (a stray word in a numeric file) is not a row;
it is counted as skipped.
\end{itemize}

\subsection{Which columns hold the attitude}\label{sec:data-roles}

wuffle then looks for columns by name. Case, spaces and punctuation don't matter: \col{Dip Direction},
\col{dip\_direction} and \col{DIPDIR} are the same name.

\begin{center}\small
\begin{tabularx}{\linewidth}{@{}>{\raggedright}p{0.24\linewidth}L@{}}\toprule
\sffamily\bfseries role & \sffamily\bfseries names wuffle recognizes \\\midrule
dip direction & dipdir, dip direction, dip azimuth, dipaz, ddir, dd, dipdirn, azimuth of dip \\
strike & strike, str, rhr, strike rhr, strike dir \\
dip & dip, dip angle, inclination, incl \\
trend & trend, tr, trd, lineation trend, azimuth, az, bearing \\
plunge & plunge, pl, plg, lineation plunge \\
sense (faults) & sense, slip sense, kinematics, kin, movement \\
angle (small circles) & angle, opening, semiangle, cone, half angle \\
set & set, family, group, class, domain \\
\bottomrule\end{tabularx}
\end{center}

The columns it finds decide the kind of dataset, in this order:

\begin{description}
\item[Faults.] A plane (dip direction or strike, and dip) \emph{and} a line (trend and plunge) on the
same row, with a sense column if there is one. \cref{ch:faults} is about them.
\item[Small circles.] A trend, a plunge and an angle: the axis and the half-angle of a cone.
\item[Planes.] Dip direction and dip, or strike and dip. A strike is read by the
\term[right-hand rule]{right-hand rule}: the plane dips to the right of the strike direction, so strike
\code{030} with dip \code{40} is the plane \code{120/40}.
\item[Lines.] Trend and plunge.
\end{description}

A \col{set} column is kept with the roles, so a later step can group by it.

\subsection{When nothing is named}\label{sec:data-guess}

Two fallbacks cover files without those names. Both are \emph{guesses}, and the table's heading says so
(\cref{sec:table}):

\begin{itemize}
\item \textbf{One column of attitudes.}\index{quadrant notation} A column whose values look like \code{045/30} or
\code{N30E 40SE} is read as planes, with quadrant letters honored: \code{N30E 40SE} is a strike of
N30°E dipping 40° to the southeast, the plane \code{120/40}. A file with one attitude per line and
nothing else is read this way, even though its values contain spaces.
\item \textbf{The first two numeric columns} are read as dip direction and dip. If the second column has
values over 90 and the first doesn't, they are swapped, since a dip can't pass 90.
\end{itemize}

In a one-column file the arrow decides. Attitudes written with an arrow (\code{045→30}, or
\code{045>30}) are lines, and the rest are planes. A column of only lines opens as a lines dataset; a
column that mixes both opens as two datasets, \code{<name> · planes} and \code{<name> · lines}, the way a
field log does (\cref{sec:field-logs}).

Columns that are named, but not as a recognized pair, also fall to the numeric guess. A file with
\col{azimuth} and \col{inclination} reads as dip direction and dip, and the table's heading says so:
\ui{no column names wuffle knows: azimuth / inclination read as dip direction / dip}. A file with no
header at all says \ui{no named columns} instead.

\subsection{Rows that can't be read}

A row whose attitude can't be read (a dip of \code{abc}, a dip over 90, an empty cell) stays in the
dataset as it is. It is not drawn, its attitude column in the table shows \ui{—}, and it takes no part
in statistics. Fix it in the table and it is drawn at once. A layer's properties dialog says how many
rows were read: \ui{tocher · planes · 200 of 200 rows read}.

\begin{why}
Refusing a file because one row in two hundred is odd would make you fix the file in another program
first, and lose the row's other columns on the way. Keeping the row and marking it means you see what
wuffle saw, and nothing is silently dropped.
\end{why}

\subsection{Field logs}

The compass's own log (\cref{sec:publish}) is recognized by its \col{type} and
\col{dip\_dir\_or\_trend} columns. Its planes and lines become two datasets,
\code{<name> · planes} and \code{<name> · lines}, both drawn on the active net. The log itself is
kept as it came, as an \emph{original}, with its \file{.sig} and photos if they are opened with it (\cref{sec:field-logs}).

\begin{tryit}
\begin{enumerate}
\item Make a text file \file{quadrant.txt} with three lines: \code{N30E 40SE}, \code{N30E 40SE} and
\code{S10W 30W}, and open it.
\item Click \ui{quadrant} in the tree's \ui{data} section and look at the table. The first column reads
\code{120/40}, \code{120/40}, \code{280/30}: the quadrant notation converted to dip direction and dip.
\item The heading ends with \ui{“c1” read as dip direction / dip (quadrant letters honored)}: the
file had no header, so the reading is marked as a guess.
\end{enumerate}
\end{tryit}

\section{The table}\label{sec:table}

\shot{data-table}{A file with a quoted field and an unreadable row: the attitude column shows \ui{—} for row 4,
which isn't drawn. Row 3 is selected, so it is highlighted here and violet on every net.}

Click a dataset in the tree's \ui{data} section, or a layer, and the \ui{table} tab\index{table tab}
shows its rows.

\begin{itemize}
\item The heading names the dataset, its kind and its row count, and when the column roles were guessed,
how they were read.
\item The first column is the \term{attitude} wuffle read from the row: dip direction and dip for planes
and faults (headed \ui{dip dir/dip}), trend and plunge for lines (\ui{trend→plunge}), and the axis of a
small circle (\ui{axis trend→plunge}). It is computed, not stored, and can't be edited. \ui{—} marks a row it
couldn't read.
\item The other columns are the file's, as they came.
\end{itemize}

\subsection{Editing}

Double-click a cell (or select it and press \key{F2}, or start typing), type, and press \key{Enter}.
The edit goes into the dataset, so every layer of it, on every diagram, redraws, and the attitude column
follows. A number typed is stored as a number; anything
else is kept as text. Each edit is an undo step, \ui{edit rows}.

\subsection{Selecting}

Click a row, or drag down several, and they are selected everywhere: violet on every net that draws the
dataset, counted in the status line. Rows selected on a net are washed in the table the same way.
\cref{ch:select} is about selections.

\section{The data section}\label{sec:data-section}

The foot of the layers tree lists every dataset in the project, with its row count. A dataset that no
diagram draws says \ui{· not drawn}; a derived one has an arrow instead of the table icon, and hovering
it says what made it (\cref{sec:lineage}).

You can drag a dataset from here onto a diagram's name in the tree to draw it there. Its right-click menu
has:

\begin{description}
\item[\ui{Show the table}] the table tab, with this dataset.
\item[\ui{Draw on}] a diagram to draw it on, as a new layer.
\item[\ui{Rename…}] a new name. Layers that use the dataset's name follow it.
\item[\ui{Delete dataset}] asks first, naming the diagrams it is drawn on, then removes the dataset
and every layer of it. Undo brings both back.
\end{description}

\begin{note}
In the console the same delete is stricter: \code{rm data/<name>.csv} refuses while the data is drawn
anywhere (\cref{sec:project-write}). The menu can ask you; a script can't, so it has to remove the layers
first.
\end{note}

\section{Lineage}\label{sec:lineage}

A dataset made by wuffle rather than read from a file carries its \term{lineage}: what made it, from
which datasets, with which settings and when. A mean, the eigenvectors, a $\beta$ axis, an intersection
(\cref{ch:stats}), a selection made into a dataset (\cref{ch:select}), and a field log's data all have
one. In the tree it shows as the hover text, \ui{made by selection from tocher}; in the project it is in
\file{data/<name>.json} (\cref{ch:project-fs}).

The lineage names its sources by name as well as by identity, so it still reads correctly after the
source is renamed or deleted.
